% concatenated from main.tex
\documentclass{amsart}

\usepackage{cleveref}
\usepackage{graphicx}
\usepackage{geometry}
\usepackage{amssymb}
\usepackage{amsmath}
\usepackage{amsfonts}
\usepackage{enumitem}
\usepackage{comment}
\usepackage{overpic}
\usepackage{amsaddr}

\newtheorem{theorem}{Theorem}[section]
\newtheorem{proposition}[theorem]{Proposition}
\newtheorem{corollary}{Corollary}[section]
\newtheorem{lemma}[theorem]{Lemma}
\theoremstyle{definition}
\newtheorem{definition}[theorem]{Definition}
\newtheorem{example}[theorem]{Example}
\newtheorem{xca}[theorem]{Exercise}
\theoremstyle{remark}
\newtheorem{remark}{Remark}

\usepackage{xcolor}
\definecolor{matlabblue}{rgb}{0,0.4470,0.7410}
\definecolor{matlabred}{rgb}{0.8500,0.3250,0.0980}
\definecolor{matlabyellow}{rgb}{0.9290,0.6940,0.1250}

\usepackage[obeyFinal,textsize=footnotesize]{todonotes}
\newcommand{\alex}[2][]{\todo[color=yellow, #1]{\textbf{Alex}: #2}}
\newcommand{\sung}[2][]{\todo[color=green, #1]{\textbf{Sungwoo}: #2}}

\title{Convergence of Pivoted Cholesky for Lipschitz Kernels}

\author{Sungwoo Jeong and Alex Townsend}
\email{sjeong@cornell.edu, townsend@cornell.edu}
\address{Department of Mathematics, Cornell University}
\date{}

\begin{document}

\begin{abstract}
    We investigate the continuous analogue of the Cholesky factorization, namely the pivoted Cholesky algorithm. Our analysis establishes quantitative convergence guarantees for kernels of minimal smoothness. We prove that for a symmetric positive definite Lipschitz continuous kernel $K:\Omega\times \Omega \rightarrow \mathbb{R}$ on a compact domain $\Omega\subset\mathbb{R}^d$, the residual of the Cholesky algorithm with any pivoting strategy is uniformly bounded above by a constant multiple of the fill distance of pivots. In particular, our result implies that under complete pivoting (where the maximum value of the diagonal of the residual is selected as the next pivot):
    \begin{equation*}
        \|R_n\|_{\infty} = \mathcal{O}(n^{-1/d}),
    \end{equation*}
    where $R_n$ is the residual after $n$ Cholesky steps and $\|\cdot\|_\infty$ is the absolute maximum value of $R_n$. Moreover, if $K$ is differentiable in both variables with a Lipschitz derivative, our convergence rate improves to $\mathcal{O}(n^{-2/d})$. Our result closes a gap between theory and practice as previous analyses required $C^2$-regularity of $K$ to establish convergence, whereas empirical evidence indicated robust performance even for non-differentiable kernels. We further detail how our convergence results propagate to downstream applications, including discrete analogues, Gaussian process regression, and the P-greedy interpolation method.
\end{abstract}

\maketitle

\section{Introduction}
The Cholesky decomposition is one of the cornerstone factorizations in numerical linear algebra, providing an efficient and stable way to represent a symmetric positive definite (SPD) matrix as a product of a lower-triangular matrix and its transpose. It underlies fast solvers for linear systems, preconditioners, and algorithms across statistics, optimization, and scientific computing. While the discrete case is well-understood, a natural extension to SPD kernels is studied far less. The Cholesky algorithm with a pivoting strategy constructs successive low-rank approximations of a kernel by iteratively selecting pivot points and adding rank-one updates (see~\cref{eq:CholeskyAlgorithm}). Despite the pivoted Cholesky algorithm having many downstream applications and related variants such as adaptive cross approximation (ACA)~\cite{bebendorf2000approximation}, pseudoskeleton approximation~\cite{goreinov1997theory}, and P-greedy interpolation~\cite{de2005near}, its convergence theory is not fully developed.

The pivoted Cholesky algorithm is best viewed as an iterative low-rank approximation scheme, rather than as a factorization. Formally, let $\Omega\subset\mathbb{R}^d$ be a compact domain. Given a SPD kernel\footnote{A kernel is called SPD if for any finite set of points $x_1,\ldots,x_n \in \Omega$, the matrix $\big[K(x_j,x_k)\big]_{j,k=1}^n$ is SPD.} $K:\Omega\times \Omega \rightarrow \mathbb{R}$, the first step of the pivoted Cholesky algorithm sets $R_0 = K$ and selects a point $z_1$ (which we call a pivot) based on prescribed selection criteria, which we call the \textit{pivoting strategy}. Then, the algorithm constructs a rank-$1$ approximation of $K$ as
\begin{equation*}
K_1(x,y) = \frac{R_0(x,z_1)R_0(z_1,y)}{R_0(z_1,z_1)},
\end{equation*}
which interpolates $K$ whenever $x = z_1$ or $y = z_1$ and the first residual $R_1=K - K_1$ is obtained. We iteratively apply this procedure. At the $n\geq 1$ step, we select a pivot point $z_n$ according to the pivot strategy and compute the rank-$n$ approximant of $K$ by adding a rank-$1$ kernel to a previous approximant, i.e., 
\begin{equation}
    K_n(x, y) = K_{n-1}(x, y) + \frac{R_{n-1}(x, z_n)R_{n-1}(z_n, y)}{R_{n-1}(z_n, z_n)},
    \label{eq:CholeskyAlgorithm}
\end{equation}
with the residual $R_n=K - K_n$. The approximant $K_n$ interpolates $K$ whenever $x = z_i$ or $y = z_i$, for $i\leq n$. The choice of pivoting strategy is crucial and often determines the quality of the approximation. Various pivoting strategies such as complete pivoting~\cite{golub2013matrix} (where $z_n = \arg\max_{x\in\Omega} |R_{n-1}(x,x)|$), maximum volume pivoting~\cite{gu1996efficient,hong1992rank} (where $z_n$ maximizes the absolute determinant of a certain submatrix), and rook pivoting~\cite{neal1992geometric} are commonly used in practice. Since $K$ is positive definite, $R_n$ is positive semidefinite for all $n$, $R_n(x,x)\geq 0$ for all $x\in\Omega$, and $\|R_n\|_\infty = \max_{x,y\in\Omega} |R_n(x,y)|$~\cite{townsend2015continuous}.



The central question is then: How well does $K_n$ approximate $K$ as $n$ grows? A classical theory~\cite{schmidt1907theorie} provides an answer for the best-case scenario. As proved by Schmidt in 1907~\cite{schmidt1907theorie}, the accuracy of any rank-$n$ approximation cannot be better than the $(n+1)$st singular value of $K$. Throughout, we assume that $K$ is a square-integrable kernels, so that we have
\begin{equation*}
     \operatorname{vol}(\Omega) \|R_n\|_{\infty} \geq \|R_n\|_2 \geq \sigma_{n+1}(K),
\end{equation*}
where $\sigma_{n+1}(K)$ denotes the $(n+1)$st singular value of $K$, $\|\cdot\|_2$ is the operator (spectral) norm, and $\|\cdot\|_{\infty}$ is the absolute maximum value on $\Omega\times \Omega$. For example, take the Brownian kernel $K : [0,1]^d\times [0,1]^d\rightarrow \mathbb{R}$ given by $K(x,y) = \prod_{i=1}^d\min\{x_i,y_i\}$. Then, $\sigma_{n}(K) \geq C_d \log^{2(d-1)}(n)/n^2$ for some constant $C_d$ that only depends on $d$ so $\|R_n\|_\infty$ can at best converge to zero at that rate. 

Bounding $\|R_n\|_\infty$ from above has long been studied for smooth kernels. It has been shown that $\|R_n\|_\infty$ decays exponentially fast to zero as $n\rightarrow\infty$ when $d = 1$, the kernel is more than analytic in each variable, and complete pivoting is used~\cite{cortinovis2020maximum,townsend2015continuous}. Moreover, in $d=1$, Bebendorf has an error bound on $\|R_n\|_\infty$ for so-called asymptotically smooth kernels when rook pivoting is employed, though that typically means the kernel must be infinitely differentiable to achieve convergence~\cite{bebendorf2000approximation}. Another convergence theory for translation invariant kernels, i.e., $K(x,y) = \phi(\|x-y\|)$ for some univariate function $\phi$, is implied by bounds on the power function from the radial basis function literature~\cite{santin6convergence}. Another notable result in the radial basis function literature suggests that the Cholesky algorithm with complete pivoting converges for all kernels $K:\Omega\times \Omega\rightarrow \mathbb{R}$ that are twice continuously differentiable and $\Omega$ is a convex domain~\cite[Thm.~4.3]{de2005near}. 

Unfortunately, these existing results do not explain the convergence of a Cholesky algorithm on kernels with very limited smoothness that arise throughout computational mathematics and applications. Prominent examples include:
(i) the Ornstein--Uhlenbeck (exponential) kernel $K(x,y)=\alpha e^{-|x-y|/\ell}$ with $\alpha,\ell>0$ (merely Lipschitz at $x=y$);
(ii) Matérn kernels with low regularity,\footnote{Here, $\mathbf{K}_\nu$ is the modified Bessel function, $\Gamma$ is the gamma function, and $\|\cdot\|$ denotes the Euclidean distance.}
\[
K_\nu(x,y)=\frac{2^{1-\nu}}{\Gamma(\nu)}\!\left(\frac{\sqrt{2\nu}\,\|x-y\|}{\ell}\right)^\nu
\mathbf{K}_\nu\!\left(\frac{\sqrt{2\nu}\,\|x-y\|}{\ell}\right),\qquad \nu,\ell >0,
\]
whose smoothness is controlled by $\nu$ and for small $\nu$ the kernel is non-differentiable at the diagonal;
(iii) Brownian-motion (Wiener) kernels on $[0,1]$, $K(x,y)=\min\{x,y\}$;
and (iv) Green’s functions of uniformly elliptic self-adjoint PDEs, such as
$G(x,y) = -\frac{1}{2\pi}\log\|x-y\|$ for the 2D Laplacian, $G(x,y)=\frac{1}{4\pi\|x-y\|}$ for the 3D Laplacian, or on $[0,1]$ with Dirichlet data $G(x,y)=\min\{x,y\}-xy$ for $-u''=f$. 
In practice, the Cholesky algorithm with any reasonable pivoting strategy continues to perform remarkably well for these nonsmooth kernels~\Cref{fig:Convergence}.

\begin{figure}[h]
  \centering
  \begin{minipage}{.49\textwidth}
    \centering
    \begin{overpic}[width=\textwidth]{figures/min_convergence}
    \put(60,35) {\rotatebox{-37}{$\mathcal{O}(n^{-1})$}}
    \put(-3,30) {\rotatebox{90}{$\|R_n\|_\infty$}}
    \put(50,-3) {$n$}
    \end{overpic}
  \end{minipage}
  \hfill
  \begin{minipage}{.49\textwidth}
    \centering
    \begin{overpic}[width=\textwidth]{figures/matern_convergence}
    \put(65,45) {\rotatebox{-45}{$\mathcal{O}(n^{-1})$}}
    \put(-3,30) {\rotatebox{90}{$\|R_n\|_\infty$}}
    \put(50,-3) {$n$}
    \end{overpic}
  \end{minipage}
  \caption{Convergence results for the Cholesky algorithm with complete pivoting for $K(x,y)=\min\{x+1,y+1\}$ with $\Omega = [-1,1]$ (left) and the Mat\'{e}rn kernel with $\nu = 1/2$, $\ell = 1/2$, and $\Omega = [-1,1]$.}
  \label{fig:Convergence}
\end{figure}

\subsection{Our contributions}
In this paper, we prove error bounds for the pivoted Cholesky algorithm applied to a SPD kernel with Lipschitz continuity along the diagonal, i.e., there is some uniform constant $L>0$ such that
\begin{equation}
|K(x, x) - K(x, y)| \leq L \|x-y\|, \qquad \forall x,y\in\Omega. 
\label{eq:LipschitzAlongDiagonal}
\end{equation} 
This condition is implied by Lipschitz continuous, but is far weaker as it allows the kernel to be potentially discontinuous off the diagonal. The error bound depends on the so-called fill distance of the pivots, which is the radius of the largest empty hypersphere that can fit among the pivots, giving us a notion of how covered $\Omega$ is by the chosen points. If the first $n$ pivots are $Z_n = \left\{z_1,\ldots,z_n\right\}\subset\Omega$, then the fill distance is given by 
\begin{equation}
h_{Z_n} = \sup_{x\in\Omega} \min_{z\in Z_n} \|x-z\|.
\label{eq:fillDistance}
\end{equation}

\subsection*{Contribution 1} 
We prove the following error bound on $\|R_n\|_\infty$ for the Cholesky algorithm with any pivoting strategy (see~\Cref{thm:filldistancebound}): 
\[
\|R_n\|_{\infty} \leq  4L h_{Z_n}.
\]
Remarkably,~\cref{eq:LipschitzAlongDiagonal} is only a condition on the diagonal of $K$, so the kernel is allowed to be potentially discontinuous away from the diagonal. Once a particular pivoting strategy is selected for the Cholesky algorithm, we can study how $h_{Z_n}\rightarrow0$ as $n\rightarrow \infty$ to derive a rate of convergence from~\Cref{thm:filldistancebound}. 

\subsection*{Contribution 2} 
For the Cholesky algorithm with complete pivoting, we have the following convergence rate (see~\Cref{thm:convergencegeneral}): 
\[
\|R_n\|_{\infty} \leq \frac{8LR}{n^{1/d}-1} = \mathcal{O}(n^{-1/d}).
\] 
where $R$ is the radius of a ball containing $\Omega\subset \mathbb{R}^d$. In~\Cref{sec:generalsmooth}, we extend~\Cref{thm:convergencegeneral} to include $C^{1, 1}$ kernels (Lipschitz continuous first derivatives), yielding an improved convergence rate of $\|R_n\|_\infty = \mathcal{O}(n^{-2/d})$ with complete pivoting. 


\subsection*{Contribution 3}
There are also matrix counterparts to~\Cref{thm:convergencegeneral}, which we find intriguing, giving us a decay rate for the Cholesky decomposition applied to structured matrices. 
Let $A$ be an $m\times m$ SPD matrix. Then, for $1<n\leq m$, we have (see~\Cref{cor:matrixversion})
\begin{equation*}
        \|A - A_n\|_{\max} \leq \frac{4(m-1)G_A}{n-1} = \mathcal{O}(n^{-1}), 
    \end{equation*}
    where $G_A$ is a discrete diagonal Lipschitz constant given by $G_A = \max_{i\neq j}|A_{i, i}-A_{i, j}|/|i-j|$. 

\subsection{Structure of the paper}

The remainder of the paper is organized as follows. In~\Cref{sec:proof} we present our main error bound on the Cholesky algorithm with any pivoting strategy for kernels satisfying~\cref{eq:LipschitzAlongDiagonal}, and then specialize this to complete pivoting to obtain a convergence rate in~\cref{sec:proofstep2}. In~\cref{sec:OtherPivotingStrategies}, we discuss other possible pivoting strategies. We improve the convergence rate for $C^{1,1}$ regularity and explain a barrier in our argument for further improvements in~\Cref{sec:generalsmooth}. In~\Cref{sec:connections} we explore the consequences of our results for matrices, Gaussian process regression, and P-greedy interpolation. Finally, we give brief concluding remarks in~\cref{sec:conclusion}.

\section{Error Bound on the Pivoted Cholesky Algorithm in Terms of Fill Distance}\label{sec:proof}
Initially, proving the convergence of the Cholesky algorithm with any pivoting strategy on nonsmooth kernels seems extremely challenging. Suppose one starts with a Lipschitz kernel with constant $L>0$, i.e., 
\[
|K(x_1,y_1) - K(x_2,y_2)| \leq L\left\| \begin{bmatrix}x_1\\y_1\end{bmatrix} - \begin{bmatrix}x_2\\y_2\end{bmatrix}\right\|,\qquad x_1,x_2,y_1,y_2\in\Omega. 
\] 
where $\|\cdot\|$ is the Euclidean norm in $\mathcal{R}^{2d}$. 
It is easy to convince oneself that this might be hopeless because, at a quick glance, it looks like the Lipschitz condition is not preserved by a Cholesky step. Since
\[
R_1(x,y) = K(x,y) - \frac{K(x,z_1)K(z_1,y)}{K(z_1,z_1)},
\]
where $z_1$ is the first pivot, you find that even with complete pivoting, one only has  
\[
\begin{aligned}
|R_1(x_1,y_1) - R_1(x_2,y_2)| &= \left|K(x_1,y_1) - \frac{K(x_1,z_1)K(z_1,y_1)}{K(z_1,z_1)} -K(x_2,y_2) + \frac{K(x_2,z_1)K(z_1,y_2)}{K(z_1,z_1)} \right|\\
&\leq L\left\| \begin{bmatrix}x_1\\y_1\end{bmatrix} - \begin{bmatrix}x_2\\y_2\end{bmatrix}\right\| + L\left\| \begin{bmatrix}z_1\\y_1\end{bmatrix} - \begin{bmatrix}z_1\\y_2\end{bmatrix}\right\|,
\end{aligned} 
\]
where in the last inequality we used the triangle inequality and that $|K(x_1,z_1)/K(z_1,z_1)|\leq 1$ and $|K(x_2,z_1)/K(z_1,z_1)|\leq 1$ for complete pivoting. This looks like bad news, as the Lipschitz condition is not preserved after one step. However, progress can be made for kernels that satisfy an even weaker condition, such as Lipschitz on the diagonal (see~\cref{eq:LipschitzAlongDiagonal}), by focusing on bounding the diagonal of $R_n$ near a pivot location. 

We start by proving that the diagonal of $R_1$ must grow slowly away from a pivot location $z_1$. The following lemma holds for the Cholesky algorithm with any pivoting strategy. 
\begin{lemma}\label{lem:lipschitzchange}
Let $K:\Omega\times \Omega\to\mathbb{R}$ be a continuous SPD kernel satisfying~\cref{eq:LipschitzAlongDiagonal} with constant $L>0$. Then, 
\begin{equation*}
    0\leq R_1(x, x) \leq 4L \|x - z_1\|, \qquad x\in\Omega,
\end{equation*} 
where $z_1$ is the first pivot selected by the Cholesky algorithm. 
\end{lemma}
\begin{proof}
Because $K$ is a positive definite kernel, so is $R_1$ and hence $R_1(x,x)\geq 0$ for all $x\in\Omega$.
Moreover, since $K$ satisfies~\cref{eq:LipschitzAlongDiagonal}, we have $K(x, x) \leq K(z_1, z_1) + 2L\|x-z_1\|$ for any $x\in\Omega$. Thus, we have 
\begin{equation}\label{eq:lemmadiagterm}
    R_1(x, x) = K(x, x) - \frac{K(x, z_1)^2}{K(z_1, z_1)} \leq K(z_1, z_1) + 2L\|x-z_1\| - \frac{K(x, z_1)^2}{K(z_1, z_1)}. 
\end{equation}
Furthermore, since $|K(z_1, z_1) - K(x,z_1)| \leq L \|x-z_1\|$ by~\cref{eq:LipschitzAlongDiagonal}, we have
\begin{equation*}
    K(z_1, z_1) - L\|x-z_1\| \leq K(x, z_1) \leq K(z_1, z_1) + L\|x-z_1\|. 
\end{equation*}
If $0\leq K(z_1, z_1)-L\|x-z_1\|$ then simply $-K(x, z_1)^2\leq -(K(z_1, z_1)- L\|x-z_1\|)^2$, and the right hand side of~\cref{eq:lemmadiagterm} is immediately bounded above by $4L\|x-z_1\|$ since we have
\begin{equation*}
    K(z_1, z_1) + 2L\|x-z_1\| - \frac{(K(z_1, z_1)-L\|x-z_1\|)^2}{K(z_1, z_1)} = 4L\|x-z_1\| - \frac{L^2\|x-z_1\|^2}{K(z_1, z_1)}\leq 4L \|x-z_1\|.
\end{equation*}
Otherwise, if $K(z_1, z_1)-L\|x-z_1\| < 0 $, then there exists $0<L'<L$ such that $K(z_1, z_1)-L'\|x-z_1\|=0$. In this case, the right hand side of~\cref{eq:lemmadiagterm} is bounded above by
\begin{equation*}
    K(z_1, z_1) + 2L\|x-z_1\| = (2L + L')\|x-z_1\| \leq 3L\|x-z_1\|.
\end{equation*}
Thus, in both cases, we find that $R_1(x, x) \leq 4 L \|x - z_1\|$ for any $x\in\Omega$ and the result follows. 
\end{proof}
\Cref{lem:lipschitzchange} tells us that along the diagonal, $R_1$ is not only zero at the pivot location, i.e., $R_1(z_1,z_1) = 0$, but it remains small in a neighborhood around it. In fact, along the diagonal, $R_1$ cannot grow faster than proportional to the distance away from the pivot. Surprisingly, this property continues to hold as one takes more Cholesky steps, regardless of the pivoting strategy. 
\begin{lemma}\label{lem:MultiStepLipschitzControl}
Let $K:\Omega\times \Omega\to\mathbb{R}$ be a SPD kernel satisfying~\cref{eq:LipschitzAlongDiagonal} with constant $L>0$. Then, the $n$th residual from the Cholesky algorithm with any pivoting strategy satisfies
\begin{equation*}
   0\leq R_n(x, x) \leq 4L \min_{1\leq i\leq n} \|x - z_i\|, \qquad x\in\Omega,
\end{equation*} 
where $z_i$ is the $i$th pivot.
\end{lemma} 
\begin{proof} 
Since $K$ is positive definite, $R_n$ is positive semidefinite and thus, $R_n(x,x)\geq 0$ for all $x\in\Omega$. 
Select the pivot location closest to $x$, say $z_i$. Since the residual $R_n$ constructed by a pivoted Cholesky algorithm is invariant to the order of the pivots, we obtain the same residual $R_n$ if we perform Cholesky steps with pivots in the order of $z_i, z_1, \dots, z_{i-1}, z_{i+1}, \dots, z_n$, instead of $z_1, \dots, z_n$. Let $\tilde{R}_1$ be the residual after performing a Cholesky step on $K$ with the pivot $z_i$ first. Then, by~\Cref{lem:lipschitzchange}, we have
    \begin{equation*}
 \tilde{R}_1(x, x) \leq 4L \|x-z_i\|,\qquad x\in\Omega,
    \end{equation*}
    since $\tilde{R}_1(z_i, z_i) = 0$. Performing the remaining $n-1$ Cholesky steps on $\tilde{R}_1$ with pivots $z_1, \dots, z_{i-1}$, $z_{i+1}, \dots, z_n$ we obtain $R_n$. Since the diagonal of the residual is nonincreasing under Cholesky steps, i.e., $R_{k+1}(x,x)\leq R_{k}(x,x)$ for all $x\in\Omega$, we find that $R_n(x, x)\leq \tilde{R}_1(x, x)$ for any $x\in\Omega$.
\end{proof} 

In particular,~\Cref{lem:MultiStepLipschitzControl} tells us that near any pivot location, $R_n$ is small. Or, said more precisely, for any $\epsilon>0$, we know that $R_n(x,x)\leq \epsilon$ provided $\min_{1\leq i\leq n} \|x - z_i\|\leq \epsilon/(4L)$. Once the diagonal of $R_n$ is controlled, one can bound $\|R_n\|_\infty$. We can now prove a bound on $\|R_n\|_\infty$ in terms of the fill distance of the pivots (see~\Cref{thm:filldistancebound}). 

\begin{theorem}\label{thm:filldistancebound}
    Let $K:\Omega\times\Omega \rightarrow\mathbb{R}$ be a SPD kernel such that~\cref{eq:LipschitzAlongDiagonal} holds for some constant $L>0$. Then, the $n$th residual $R_n$ of the pivoted Cholesky algorithm satisfies
\[
        \|R_n\|_{\infty} \leq  4L h_{Z_n},
\]
    where $Z_n=\{z_1, \dots, z_n\}$ is the set of pivots and $h_{Z_n}$ is defined in~\cref{eq:fillDistance}. 
\end{theorem}
\begin{proof}
Let $Z_n = \{z_1,\ldots,z_n\}$ be the pivots selected by Cholesky. Since $R_n(z_i,z_i) = 0$ for $1\leq i\leq n$, we know from~\Cref{lem:MultiStepLipschitzControl} that 
\[
R_n(x,x)\leq 4L \min_{z_i\in Z_n} \|x - z_i\|\leq 4L h_{Z_n}, \qquad x\in\Omega, 
\]
where the last inequality follows directly from the definition of $h_{Z_n}$ (see~\cref{eq:fillDistance}).  Finally, since $R_n$ is itself a SPD kernel, we know that\footnote{This follows because the $2\times 2$ matrix $\begin{bmatrix} R(x,x) & R(x,y) \\ R(y,x) & R(y,y)\end{bmatrix}$ is positive definite.}
\[
|R_n(x,y)| \leq \sqrt{R_n(x,x)R_n(y,y)} \leq 4L h_{Z_n}, \qquad x,y\in\Omega,
\]
as desired.
\end{proof}

Just because $\|R_n\|_\infty\leq 4Lh_{Z_n}$ does not necessarily imply that $\|R_n\|_{\infty}$ to 0 as $n\to\infty$, since the fill distance $h_{Z_n}$ need not shrink.   In fact, it may be plausible that the pivot points could accumulate in such a way that the fill distance is not reduced as $n\rightarrow \infty$. One pivoting strategy one could employ to minimize the fill distance is uniform pivoting, where the pivot locations are predetermined and selected based on $\Omega$, not $K$. One can select a set of (quasi-)uniform points in $\Omega$ as the pivot locations. If $\Omega=[0,1]^d$ and $Z_n$ is taken as a tensor grid with $m$ points per coordinate (so $n=m^d$), then every $x\in[0,1]^d$ lies within Euclidean distance at most $\frac{\sqrt d}{2m}$ of some grid point. Hence
\[
h_{Z_n}\;\le\;\frac{\sqrt d}{2m}\;=\;\frac{\sqrt d}{2}\,n^{-1/d},
\]
and by \Cref{thm:filldistancebound} we obtain the explicit bound
\begin{equation*}
    \|R_n\|_{\infty}\;\le\;4L\,h_{Z_n}\;\le\;2L\,\sqrt d\,n^{-1/d}.
\end{equation*}
More generally, for any quasi-uniform set $Z_n$ with bounded mesh ratio, one has $h_{Z_n}\le C_d\,n^{-1/d}$, yielding $\|R_n\|_{\infty}\le 4L\,C_d\,n^{-1/d}$ with a constant $C_d$ depending only on $d$ (and the mesh ratio). 

The downside to uniform pivoting is that it does not exploit the properties of $K$, and its empirical convergence rate does not always closely match the decay rate of the kernel's singular values. In the next section, we show that the Cholesky algorithm with complete pivoting also ensures that the fill distance goes to $0$ while also converging much faster in practice.

\section{On the convergence of the Cholesky algorithm with complete pivoting}\label{sec:proofstep2}
Now, we focus on the Cholesky algorithm with complete pivoting. Since complete pivoting always selects the next pivot as the absolute maximum of the residual, one can show that two pivots cannot be close unless the residual is already small. (To get a sense of the pivot locations selected by the Cholesky algorithm with complete pivoting, see~\Cref{fig:pretty}.)

\begin{figure} 
\begin{minipage}{.32\textwidth}
\includegraphics[width=\textwidth]{figures/pretty1.eps}
\end{minipage} 
\begin{minipage}{.32\textwidth}
\includegraphics[width=\textwidth]{figures/pretty2.eps}
\end{minipage} 
\begin{minipage}{.32\textwidth}
\includegraphics[width=\textwidth]{figures/pretty3.eps}
\end{minipage} 
\begin{minipage}{.32\textwidth}
\includegraphics[width=\textwidth]{figures/pretty4.eps}
\end{minipage} 
\begin{minipage}{.32\textwidth}
\includegraphics[width=\textwidth]{figures/pretty5.eps}
\end{minipage} 
\begin{minipage}{.32\textwidth}
\includegraphics[width=\textwidth]{figures/pretty6.eps}
\end{minipage} 
\caption{Six contour plots of kernels with $\Omega = [-1,1]$, showing the pivot location (black dots) of the Cholesky algorithm with complete pivoting until residual is below machine precision. (Top, left) $K(x,y) = 1/(1+100(x^2-y^2)^2)$ (Top, middle) $K(x,y) = e^{-5(x-y)^2}$, (Top, right) Mat\'{e}rn kernel with $\nu = 1/2$ and $\ell =1/2$ (only showing first 10 pivots), (Bottom, left) $K(x,y) = \min\{x+1,y+1\}$ (only showing first 10 pivots), (Bottom, middle) Kernel from~\cite[Fig.~4.10 (right)]{townsend2014computing}, and (Bottom, right) $K(x,y) = 1/(1 + (x^2 + y^2))$.}
\label{fig:pretty}
\end{figure}



\begin{lemma}\label{prop:pivotdistance}
    Let $K:\Omega\times \Omega\to\mathbb{R}$ be a SPD kernel satisfying~\cref{eq:LipschitzAlongDiagonal}. The residual of the Cholesky algorithm with complete pivoting satisfies 
    \[
    \|R_n\|_\infty \leq 4L \min_{1\leq i,j\leq n} \|z_i - z_j\|,
    \]
    where $z_1,\ldots,z_n$ are the Cholesky pivots. 
\end{lemma}
\begin{proof}
    Let $z_i$ and $z_j$ be the closest two pivots with $i<j$. From~\Cref{lem:MultiStepLipschitzControl} with $n = j-1$, we see that 
    \[
    \|R_{j-1}\|_\infty = |R_{j-1}(z_j,z_j)| \leq 4L \|z_j - z_i\|
    \]
    where the first equality comes from the fact that the pivots are selected according to complete pivoting. The result follows because $\|R_0\|_\infty \geq \|R_1\|_\infty \geq \|R_2\|_\infty\geq \cdots$.
\end{proof}

Given~\Cref{prop:pivotdistance} and a given number of pivots $n$, all we have to do is find a number $r$ such that there are always two pivots $z_i, z_j$ with $i, j\leq n$ such that $\|z_i - z_j\| \leq r$. 

\begin{theorem}\label{thm:convergencegeneral}
    Let $\Omega \subset \mathbb{R}^d$ be a compact domain contained in a ball of radius $R$. Let $K:\Omega\times \Omega\to\mathbb{R}$ be a kernel that satisfies the assumptions of~\Cref{thm:filldistancebound}. Then, the $n$th residual $R_n$ of the Cholesky algorithm with complete pivoting satisfies $(n>1)$
    \begin{equation}\label{eq:generalbound}
        \|R_n\|_{\infty} \leq \frac{8LR}{n^{1/d}-1} = \mathcal{O}(n^{-1/d}). 
    \end{equation}
\end{theorem}
\begin{proof}
    For $n=1$, from \Cref{lem:MultiStepLipschitzControl}, we directly have $\|R_1\|_\infty \leq 4L \max_{x, y\in\Omega}\|x-y\| \leq 8LR$. For $n>1$, let $r = 2R/(n^{1/d}-1)$. This $r$ is designed to satisfy the following identity exactly:
    \begin{equation*}
        n = \left(1 + \frac{2R}{r}\right)^d = \frac{\operatorname{vol}(B_{R + r/2})}{\operatorname{vol}(B_{r/2})}, 
    \end{equation*}
    where $B_r$ is the closed ball in $\mathbb{R}^d$ of radius $r$ and $\operatorname{vol}$ is its volume. Denote the first $n$ pivots of Cholesky as $z_1, \dots, z_n$ and consider $n$ closed balls of radius $\frac{r}{2}$ centered at $z_1, \dots, z_n$. It is impossible for all of these $n$ balls to be disjoint, as otherwise, 
        \begin{equation*}
        n \operatorname{vol}(B_{r/2}) < \operatorname{vol}(B_{R+\frac{r}{2}}),
    \end{equation*}
    which contradicts the fact that $n = \left(1 + \frac{2R}{r}\right)^d$. Therefore, we must have $\min_{1\leq i,j\leq n}\|z_i - z_j\|\leq r$. The bound on $\|R_n\|_\infty$ in~\cref{eq:generalbound} follows from \Cref{prop:pivotdistance}.
\end{proof}

This result shows that there is a general convergence guarantee for the Cholesky algorithm with complete pivoting when applied to kernels that are only Lipschitz along the diagonal. In particular, \Cref{thm:convergencegeneral} shows that complete pivoting necessarily drives the residual $\|R_n\|_\infty$ to zero at a dimension-dependent algebraic rate $\mathcal{O}(n^{-1/d})$, regardless of the geometry of $\Omega$. This confirms that the pivots cannot accumulate, and the pivot set generated by complete pivoting always has shrinking fill distance. So, even under very weak regularity assumptions on $K$, the Cholesky algorithm with complete pivoting provides a systematic approximation procedure with provable convergence. In particular, implicit in the proof of~\Cref{thm:convergencegeneral} is the fact that the first $n$ pivots selected by the Cholesky algorithm with complete pivoting on kernels with a diagonal Lipschitz condition (see~\cref{eq:LipschitzAlongDiagonal}) satisfy 
\[
h_{Z_n} \leq \frac{2R}{n^{1/d} - 1}.
\]
where $\Omega$ is contained in a ball of radius $R$. In~\Cref{fig:ConvergenceDiffD}, we demonstrate this convergence rate for the Cholesky algorithm with complete pivoting for $d = 1, 2, 3$ with the Mat\'{e}rn kernel with $\nu=1/2$, $\ell = 1/2$, and $\Omega = [-1,1]^d$. 

\begin{figure} 
\begin{minipage}{.49\textwidth}
\begin{overpic}[width=\textwidth]{figures/convergence_diff_d}
 \put(70,47) {\rotatebox{-45}{$\mathcal{O}(n^{-1})$}}
  \put(75,67) {\rotatebox{-23}{$\mathcal{O}(n^{-1/2})$}}
   \put(80,75) {\rotatebox{-16}{$\mathcal{O}(n^{-1/3})$}}
    \put(-5,35) {\rotatebox{90}{$\|R_n\|_\infty$}}
  \put(60,35) {\color{matlabblue}\rotatebox{-45}{$d = 1$}}
  \put(75,57) {\color{matlabred}\rotatebox{-23}{$d = 2$}}
   \put(61,72.5) {\color{matlabyellow}\rotatebox{-17}{$d = 3$}}
    \put(50,-3) {$n$}
\end{overpic}
\end{minipage}
\caption{The convergence of the Cholesky algorithm with complete pivoting for the Mat\'{e}rn kernel with $\nu = 1/2$, $\ell = 1/2$, and $\Omega = [-1,1]^d$ for $d = 1$ (blue), $d=2$ (red), and $d=3$ (yellow). The convergence rate matches the $\mathcal{O}(n^{-1/d})$ rate predicted by~\Cref{thm:convergencegeneral}.}
\label{fig:ConvergenceDiffD}
\end{figure} 
\section{Other pivoting strategies}\label{sec:OtherPivotingStrategies}
The Cholesky algorithm can be employed with many different pivoting strategies. In this section, we briefly consider other pivoting strategies such as: (1) local maximum volume pivoting (see~\cref{sec:LocalMaxVol}), (2) complete pivoting with mistakes (see~\cref{sec:completePivotingWithMistakes}), and (3) discrete pivoting (see~\cref{sec:discretepivoting}).

\subsection{Global and local maximum volume pivoting}\label{sec:LocalMaxVol}

Global and local maximum volume pivoting are classical strategies in numerical linear algebra for selecting pivots in rank-revealing factorizations~\cite{damle2024reveal,gu1996efficient}. Given a set of previously chosen pivots, the idea is to select the next pivot to maximize the determinant (i.e., the volume) of the corresponding pivot submatrix. That is, the first $k$ pivots $Z_{k} = \left\{z_1,\ldots,z_k\right\}$ for global maximum volume pivoting are chosen so that
   \[
    \det\bigl(K[Z_k,Z_k]\bigr) \geq \det\bigl(K[B_k,B_k]\bigr)
    \]
    where $K[Z_k,Z_k]$ is the $k\times k$ submatrix of $K$ obtained by evaluating the kernel at $Z_k$ forming a matrix with entries $K(z_i,z_j)$ and $B_k$ is any collection of $k$ points in $\Omega$. Despite its beautiful theoretical properties, global maximum volume is usually not computationally feasible~\cite{civril2009selecting}. Local maximum volume is a more feasible alternative, where the first $k$ pivots $Z_{k+1}$ are selected so that 
    \[
    \det\bigl(K[Z_{k},Z_{k}]\bigr) \;\geq\; \det\bigl(K[\left(Z_{k}\setminus z\right)\cup\{y\},\left(Z_{k}\setminus z\right)\cup\{y\}]\bigr),
    \]
    for all $z\in Z_{k}$ and all $y\in \Omega \setminus Z_k$. That is, the next pivot is chosen so that replacing any one pivot with the new candidate would not increase the volume of the selected submatrix. Both global and local maximum volume pivoting are non-greedy pivot selection criteria. The same error bound as in~\Cref{thm:convergencegeneral} holds for the Cholesky algorithm with global or local maximum volume pivoting. 

\begin{corollary}
    With the same assumptions as \Cref{thm:filldistancebound}, the $n$th residual $R_n$ of the Cholesky algorithm with global or local maximum volume pivoting satisfies
    \begin{equation*}
        \|R_n\|_\infty \leq \frac{8LR}{n^{1/d}-1}.
    \end{equation*}
    \label{cor:maxvol}
\end{corollary}
\begin{proof}
    We will once more prove that the residual satisfies the bound $\|R_n\|_\infty \leq 4L \min_{i, j}\|z_i - z_j\|$. For a given $n$, let $z_{\ell}, z_k$ be the two pivots with the minimum distance. Consider the residual $\tilde{R}_{n-1}$ obtained by performing $n-1$ Cholesky steps with pivots $z_1, \dots, z_{k-1},z_{k+1},\ldots, z_n$. Define the set $S_k = \{1, \dots, n\}\setminus\{k\}$. The diagonal values of $\tilde{R}_{n-1}$ are given by    \begin{equation*}
        \tilde{R}_{n-1}(x, x) = K(x, x) - \left([K(x, z_i)]_{i=S_k}\right)^T [K(z_i, z_j)]_{i, j\in S_k} [K(x, z_i)]_{i=S_k},
    \end{equation*}
    and notice that 
    \begin{equation*}
        \det\left(K[S_k \cup \{y\}, S_k \cup \{y\}]\right) = \det\left(K[S_k,S_k]\right) \cdot \tilde{R}_{n-1}(y, y),
    \end{equation*}
    for all $y\in \Omega$, from the block determinant formula. If there is $y\in\Omega$ such that $\tilde{R}_{n-1}(y, y)>\tilde{R}_{n-1}(z_k, z_k)$, then the volume corresponding to pivots $S_k \cup \{y\}$ is larger than $S_k\cup\{z_k\} = \{z_1, \dots, z_n\}$, which is a contradiction as it would violate the global or local maximum volume criterion. Thus, the diagonal of $\tilde{R}_{n-1}$ must attain its maximum at $z_k$. By \Cref{lem:MultiStepLipschitzControl} and the fact that $z_{\ell}$ is the nearest pivot to $z_k$, we have that $R_n(x, x) \leq \tilde{R}_{n-1}(x, x) \leq \tilde{R}_{n-1}(z_k, z_k)\leq 4L \|z_\ell - z_k\|$. 
\end{proof}
This corollary demonstrates that maximum volume pivoting enforces the same geometric spacing of pivots as complete pivoting, and thus achieves the same convergence rate. 
However, it is essential to note that this bound is not rank-revealing, which one might hope for. While~\Cref{cor:maxvol} controls the decay of the residual uniformly, it does not guarantee that the pivot set exposes the numerical rank of the kernel matrix in the sense of strong rank-revealing factorizations~\cite{damle2024reveal}. An alternative analysis is needed to address rank-revealing properties of maximum volume pivoting.


\subsection{Complete pivoting with mistakes}\label{sec:completePivotingWithMistakes}
In software implementations, complete pivoting may not be practical because of the requirement to compute the maximum of the diagonal of the residual at every step. To resolve such a problem, one can approximately select the maximum at each step, without significantly degrading the computed approximant.
Let $0<\delta\leq 1$ be a parameter measuring the quality of a pivot. We say that $\delta$-complete pivoting is the scheme where at the $i$th Cholesky step, one selects any pivot point $z_i$ such that 
\[
R_{i-1}(z_i, z_i) \geq \delta \| R_{i-1}\|_{\infty}.
\]
For instance, when $\delta=1/10$, any point $x\in\Omega$ may be chosen as a pivot provided that $R_{i-1}(x,x)$ is at least one-tenth of the maximum diagonal value. This relaxation of complete pivoting makes finding pivots significantly easier. Moreover,~\Cref{thm:convergencegeneral} can be easily extended to $\delta$-complete pivoting to obtain: 

\begin{corollary}
    Let $0< \delta \leq 1$. With the same assumptions as~\Cref{thm:filldistancebound}, the $n$th residual $R_n$ of the Cholesky algorithm with $\delta$-complete pivoting satisfies
    \begin{equation*}
        \|R_n\|_{\infty} \leq \frac{8 LR}{\delta(n^{1/d}-1)}.
    \end{equation*}
\end{corollary}

\subsection{Discrete pivoting}\label{sec:discretepivoting}

In practice, exact complete pivoting is often impractical: at each step, one would have to evaluate the diagonal of the residual everywhere in $\Omega$ and locate its global maximum. A standard remedy is to discretize the diagonal search: fix an $\eta$-net (grid) $G\subset\Omega$ and, at each step, pick the point in $G$ with the largest diagonal residual.

For kernels that are Lipschitz on the diagonal, our one-sided control gives (see~\Cref{lem:MultiStepLipschitzControl})
\[
0\le R_{i-1}(x,x)\;\le\;4L\,\mathrm{dist}(x,Z_{i-1}),
\]
where $\mathrm{dist}(x,Z_{i-1})$ is the minimum distance between $x$ and a pivot location in $Z_{i-1}$. However, this bound is not enough to quantify how well a grid captures large residuals. Let $M=\|R_{i-1}\|_\infty$ and let $x_\star$ be a maximizer, i.e., $R_{i-1}(x_\star,x_\star)=M$. If $y\in G$ is any grid point with $\|x_\star-y\|\le\eta$, then
\[
\mathrm{dist}(x_\star,Z_{i-1})
\;\le\;\|x_\star-y\|+\mathrm{dist}(y,Z_{i-1})
\;\le\;\eta+\mathrm{dist}(y,Z_{i-1}).
\]
Multiplying by $4L$ and using $M\le 4L\,\mathrm{dist}(x_\star,Z_{i-1})$ and $R_{i-1}(y,y)\le 4L\,\mathrm{dist}(y,Z_{i-1})$ gives the additive capture bound
\[
M \;\le\; 4L\eta \;+\; R_{i-1}(y,y).
\]
Hence, the grid maximizer $y_G\in G$ satisfies
\[
R_{i-1}(y_G,y_G)\;\ge\; M-4L\eta.
\]
In particular, whenever $4L\eta\le (1-\delta)M$, the grid pick is a \(\delta\)-approximate pivot:
\[
R_{i-1}(y_G,y_G)\;\ge\;\delta\,\|R_{i-1}\|_\infty.
\]
Thus, discretizing the search reduces exact complete pivoting to $\delta$-complete pivoting with a tolerance that is explicitly controlled by the grid spacing~$\eta$. As the algorithm progresses and the maximum residual $M$ shrinks, one may refine the grid (decrease $\eta$) to maintain a target~$\delta$ bounded away from~$0$.

If at every step we select a pivot satisfying 
\(
R_{i-1}(z_i,z_i)\ge \delta\,\|R_{i-1}\|_\infty
\)
(with any $0<\delta\le 1$), we find that 
\(
\|R_{j-1}\|_\infty \le (4L/\delta)\,\min_{i<j}\|z_i-z_j\|
\),
and therefore
\[
\|R_n\|_\infty \;\le\; \frac{8 L R}{\delta\,(n^{1/d}-1)}.
\]
Thus, choosing a grid with spacing $\eta$ such that $4L\eta\le (1-\delta)\|R_{i-1}\|_\infty$ guarantees that the grid-based pivot achieves at most a $\delta$ pivoting mistake at step $i$. In practice, a simple strategy is to start with a coarse grid and refine it adaptively so that $\delta$ remains near~$1$ while keeping evaluations inexpensive.  

There is an adaptive pivoting strategy that one can use for Gaussian elimination on general bivariate functions, and one can employ a similar discrete pivoting strategy for SPD kernels, see~\cite{townsend2013extension}. 

\section{Error bounds for the pivoted Cholesky algorithm on kernels with extra smoothness}\label{sec:generalsmooth}

We now turn to kernels with additional smoothness and show that stronger convergence rates can be obtained when additional smoothness is assumed on the kernel. We write $\partial_j^{1}K(x,y)$ for the $j$th component of the gradient of $K(x,y)$ with respect to its first variable, and similarly $\partial_j^{2}K(x,y)$ for derivatives with respect to the second variable. For a scalar function $f:\Omega\to\mathbb{R}$, we use $\partial_j f$ to denote the $j$th component of its gradient.

\begin{lemma}\label{lem:lipschitzchangemulti}
Let $\Omega\subset\mathbb{R}^d$ be compact and let $K:\Omega\times\Omega\to\mathbb{R}$ be a continuous SPD kernel.
Assume:
\begin{enumerate}[label=(\alph*)]
\item There exists $L>0$ such that, for all $1\le j\le d$,
\begin{equation}\label{eq:diag-deriv-lip}
\bigl|\partial_j^{1}K(x,y)-\partial_j^{1}K(x,x)\bigr|\le L\|x-y\|,
\qquad
\bigl|\partial_j^{2}K(x,y)-\partial_j^{2}K(x,x)\bigr|\le L\|x-y\|,
\quad x,y\in\Omega,
\end{equation}
where the partials exist on $\Omega\times\Omega$.
\item There exists $L_0>0$ such that
\begin{equation}\label{eq:first-arg-lip}
\bigl|K(x,y)-K(x',y)\bigr|\le L_0\|x-x'\|\qquad \text{for all }x,x',y\in\Omega,
\end{equation}
and $K$ is differentiable in its first argument on $\Omega\times\Omega$.
\end{enumerate}
Let $R_1$ be the residual after one Cholesky step (with any pivoting strategy) and let $z_1$ be the first pivot. Then the diagonal of the residual $D_1(x)=R_1(x,x)$ is differentiable on $\Omega$, and for each $j=1,\dots,d$ there exists a constant $C>0$ (depending only on $L$, $L_0$, and $K_{\min}=\min_{x\in\Omega}K(x,x)>0$) such that
\begin{equation}\label{eq:grad-D1-linear}
\bigl|\partial_j D_1(x)\bigr|\;\le\; \left(4L + \frac{2L_0^2}{K(z_1,z_1)}\right)\,\|x-z_1\|\qquad \text{for all }x\in\Omega.
\end{equation}
\end{lemma}
\begin{proof}
Write
\[
D_1(x)=K(x,x)-\frac{K(x,z_1)^2}{K(z_1,z_1)},\qquad
\partial_j D_1(x)
= \partial_j^1K(x,x)+\partial_j^2K(x,x)
-\frac{2}{K(z_1,z_1)}\,K(x,z_1)\,\partial_j^1K(x,z_1).
\]
Because $R_1(\cdot,z_1)=R_1(z_1,\cdot)=0$, we have $\partial_j D_1(z_1)=0$ for all $j$.
Subtracting $\partial_j D_1(z_1)$ and using the triangle inequality,
\begin{align*}
|\partial_j D_1(x)|
&= \bigl|\partial_j D_1(x)-\partial_j D_1(z_1)\bigr| \\
&\le \underbrace{\bigl|\partial_j^1K(x,x)-\partial_j^1K(x,z_1)\bigr|}_{\le\, L\|x-z_1\|\ \text{by }\eqref{eq:diag-deriv-lip}}
+\underbrace{\bigl|\partial_j^2K(x,x)-\partial_j^2K(x,z_1)\bigr|}_{\le\, L\|x-z_1\|\ \text{by }\eqref{eq:diag-deriv-lip}} \\
&\quad + \frac{2}{K(z_1,z_1)}
\underbrace{\bigl|K(x,z_1)-K(z_1,z_1)\bigr|}_{\le\, L_0\|x-z_1\|\ \text{by }\eqref{eq:first-arg-lip}}
\ \underbrace{\bigl|\partial_j^1K(x,z_1)\bigr|}_{\le\, L_0\ \text{since }K\text{ is }L_0\text{-Lipschitz in }x} \\
&\quad + \frac{2}{K(z_1,z_1)}\,K(z_1,z_1)\,
\underbrace{\bigl|\partial_j^1K(x,z_1)-\partial_j^1K(z_1,z_1)\bigr|}_{\le\, L\|x-z_1\|\ \text{by }\eqref{eq:diag-deriv-lip}}.
\end{align*}
Here we used that \cref{eq:first-arg-lip} and differentiability in $x$ imply
\(
\sup_{(x,y)\in\Omega^2}|\partial_j^{1}K(x,y)|\le L_0
\)
(by the mean value theorem). Combining the bounds gives
\[
|\partial_j D_1(x)| \;\le\; \left(4L + \frac{2L_0^2}{K(z_1,z_1)}\right)\,\|x-z_1\|,
\]
which proves \cref{eq:grad-D1-linear}.
\end{proof}

The assumptions of~\Cref{lem:lipschitzchangemulti} are automatically satisfied by $C^{1,1}$ kernels.  Indeed, if $K\in C^{1,1}(\Omega\times\Omega)$, then all first partial derivatives of $K$ are globally Lipschitz on the compact set $\Omega\times\Omega$.  This immediately implies condition~\cref{eq:diag-deriv-lip}, since fixing $x$ and comparing $\partial_j^1 K(x,y)$ with $\partial_j^1 K(x,x)$ amounts to moving along a line of length $\|x-y\|$ in the $(x,y)$ variables.  Likewise, continuity and boundedness of $\nabla_x K$ on $\Omega\times\Omega$ ensure condition~\cref{eq:first-arg-lip} via the mean value theorem, with constant $L_0=\sup_{(x,y)\in\Omega^2}\|\nabla_x K(x,y)\|$. 
Thus, every $C^{1,1}$ kernel satisfies both (a) and (b), so the lemma applies in particular to this class of kernels. 

\begin{lemma}\label{lem:boundedregiondiff2}
Assume the hypotheses of \Cref{lem:lipschitzchangemulti} and $K_{\min} = \min_{z\in\Omega} K(z,z)>0$. 
Let $z_1,\dots,z_n\in\Omega$ be the first $n$ Cholesky pivots (for any pivoting strategy). Then,
\begin{equation}\label{eq:quadratic-near-pivots}
0\;\le\; R_n(x,x)\;\le\; \sqrt{d}\left(2L + \frac{L_0^2}{K_{\min}}\right)\,\min_{1\leq i\leq n} \|x - z_i\|^2, \qquad x\in\Omega.
\end{equation}
In particular, the diagonal residual is quadratic in the distance to the nearest pivot.
\end{lemma}
\begin{proof}
Fix $x\in\Omega$ and choose an index $i$ so that $\|x-z_i\|$ is minimized. 
By the order-invariance of the residual with respect to the set of chosen pivots, we may perform the first Cholesky step at $z_i$ and then apply the remaining $n-1$ steps in any order to obtain the same $R_n$. Let $\tilde{R}_1$ denote the residual after the single step at pivot $z_i$, and set $\tilde{D}_1(x)=\tilde{R}_1(x,x)$.  By~\Cref{lem:lipschitzchangemulti}, each partial derivative of $\tilde{D}_1$ satisfies the pointwise bound
\[
|\partial_j \tilde{D}_1(x)| \;\le\; \underbrace{\left(4L + \frac{2L_0^2}{K(z_i,z_i)}\right)}_{=C_1}\,\|x-z_i\|.
\]
Hence, $\|\nabla \tilde{D}_1(x)\|\le \sqrt{d}\,C_1\,\|x-z_i\|$. 
Using $\tilde{D}_1(z_i)=0$ and $\nabla \tilde{D}_1(z_i)=0$, integrate along the line segment $\gamma(t)=z_i+t(x-z_i)$, we find that
\[
\tilde{D}_1(x)
=\!\! \int_0^1 \!\!\!\nabla \tilde{D}_1(\gamma(t))\cdot (x-z_i)\,dt
\le \!\!\int_0^1 \!\!\!\|\nabla \tilde{D}_1(\gamma(t))\|\,\|x-z_i\|\,dt
\le\!\!\int_0^1 \!\!\!\sqrt{d}\,C_1\,t\,\|x-z_i\|^2\,dt
=\frac{\sqrt{d}\,C_1}{2}\,\|x-z_i\|^2.
\]
Since subsequent Cholesky steps can only decrease the diagonal of the residual,
\[
0 \le R_n(x,x) \le \tilde{D}_1(x) \le \frac{\sqrt{d}\,C_1}{2}\,\|x-z_i\|^2.
\]
The argument can be repeated for any $x\in\Omega$ to reach the conclusion in~\cref{eq:quadratic-near-pivots}. 
\end{proof}

In particular,~\Cref{lem:boundedregiondiff2} shows that for $C^{1,1}$ kernels the diagonal of the residual is locally controlled at second order near the pivot set. For $C^{1,1}$ kernels, we see that the residual vanishes at each pivot and cannot grow faster than quadratically with the distance to the nearest pivot. 
This quadratic control is a marked improvement over the merely linear behavior available under diagonal Lipschitz continuity. We can write the bound in terms of the fill distance: 
\begin{theorem}\label{thm:filldistance-squared}
Assume the hypotheses of \Cref{lem:lipschitzchangemulti} and $K_{\min} = \min_{z\in\Omega} K(z,z)>0$. Let $Z_n=\{z_1,\dots,z_n\}\subset\Omega$ be the first $n$ pivots (for any pivoting strategy). Then, for all $x\in\Omega$,
\[
0\;\le\; R_n(x,x)\;\le\; \sqrt{d}\left(2L + \frac{L_0^2}{K_{\min}}\right) h^2_{Z_n},
\]
where $h_{Z_n}$ is the fill distance of $\{z_1,\dots,z_n\}$. 
\end{theorem}
\begin{proof}
The result is immediate from~\Cref{lem:boundedregiondiff2}. 
\end{proof} 

The fact that we can bound $R_n$ on the diagonal quadratically with the distance to the nearest pivot is the key mechanism behind an improved convergence rate for both uniform pivoting and complete pivoting. 

For uniform pivoting, the quadratic diagonal control in \Cref{thm:filldistance-squared} leads to a substantially faster convergence rate. 
As before, suppose $\Omega=[0,1]^d$ and take $Z_n$ to be a tensor grid with $m$ points per coordinate (so $n=m^d$). Then, we have $h_{Z_n}\leq \frac{\sqrt d}{2}\,n^{-1/d}$ and hence,
\[
\|R_n\|_{\infty}\;\le\;\sqrt{d}\!\left(2L+\frac{L_0^2}{K_{\min}}\right)\,h_{Z_n}^2
\;\le\;C\,d^{3/2}\,n^{-2/d},
\]
for some constant $C>0$. Thus, under $C^{1,1}$ regularity of the kernel $K$, uniform pivoting improves the algebraic rate of convergence $\mathcal{O}(n^{-2/d})$ as $n\rightarrow \infty$.

For complete pivoting, the quadratic diagonal control again yields an improved rate. Since we know that complete pivoting ensures that $h_{Z_n} \leq (2R)/(n^{1/d} - 1)$, we find that 
\[
\|R_n\|_{\infty}\;\le\;\sqrt{d}\left(2L + \frac{L_0^2}{K_{\min}}\right) \left(\frac{2R}{n^{1/d} - 1}\right)^2\;=\;\mathcal{O}(n^{-2/d}),
\]
where $\Omega$ is contained in a ball of radius $R$. Thus, under $C^{1,1}$ regularity of the kernel, complete pivoting automatically enforces quasi-uniform pivot distributions and achieves the same $\mathcal{O}(n^{-2/d})$ rate as uniform pivoting, but without requiring predetermined pivoting locations. In~\Cref{fig:smoother}, we show that as the kernel gets smoother, we witness higher rates of convergence of the Cholesky algorithm with complete pivoting. 

\begin{figure}
\begin{minipage}{.49\textwidth}
\begin{overpic}[width=\textwidth]{figures/smooth_convergence}
  \put(83,28) {\rotatebox{-45}{$\mathcal{O}(n^{-3})$}}
  \put(75,50) {\rotatebox{-30}{$\mathcal{O}(n^{-2})$}}
  \put(75,65) {\rotatebox{-16}{$\mathcal{O}(n^{-1})$}}
  \put(80,19) {\color{matlabyellow}{\rotatebox{-45}{$\nu=3/2$}}}
  \put(75,40) {\color{matlabred}{\rotatebox{-30}{$\nu=1$}}}
  \put(75,56) {\color{matlabblue}{\rotatebox{-16}{$\nu=1/2$}}}
  \put(-5,35) {\rotatebox{90}{$\|R_n\|_\infty$}}
  \put(50,-3) {$n$}
\end{overpic}
\end{minipage}
\caption{The convergence of the Cholesky algorithm with complete pivoting for the Mat\'{e}rn kernel with $\ell=1/2$, $\Omega = [-1,1]$, and $d = 1$ for $\nu = 1/2$ (blue), $\nu=1$ (red), and $\nu = 3/2$ (yellow). The theory provided in this paper does not explain the $\mathcal{O}(n^{-3})$ convergence rate when $\nu=3/2$. However, since the Mat\'{e}rn kernel is translation-invariant and for $\nu\geq 3/2$ the kernel is $C^2$ smooth the P-greedy literature is a useful resource (see~\Cref{sec:Pgreedy} and~\cite{santin6convergence}).}
\label{fig:smoother}
\end{figure} 

Finally, we conclude this section with a remark that a similar argument for kernels that are several times differentiable fails to yield a higher-order convergence rate. The following example illustrates why. Let $\Omega=[0,1]$ and consider the Gaussian kernel 
\[
K(x,y) = \exp\!\Big(-\frac{(x-y)^2}{2\sigma^2}\Big), \qquad \sigma>0,
\]
which is infinitely many times differentiable. Fix a pivot $z\in[0,1]$ and perform a single Cholesky step at $z$. On the diagonal, the residual is
\[
R_1(x,x) \;=\; 1 - \frac{K(x,z)^2}{K(z,z)}
\;=\; 1 - \exp\!\Big(-\frac{(x-z)^2}{\sigma^2}\Big).
\]
A Taylor expansion of the exponential at $x=z$ gives
\[
\frac{(x-z)^2}{\sigma^2} - \frac{(x-z)^4}{2\sigma^4}
\;\le\; R_1(x,x) \;\le\; \frac{(x-z)^2}{\sigma^2},
\]
so there exist constants $c_\sigma,r_\sigma>0$ such that
\[
R_1(x,x) \;\ge\; c_\sigma\,|x-z|^2 \qquad\text{whenever } |x-z|\le r_\sigma.
\]
Thus, even with extra regularity, we only have a quadratic bound on the residual around pivot points. 

\section{Consequences of our convergence rates}\label{sec:connections}

The pivoted Cholesky algorithm appears in a number of applications. However, in most applications, the theoretical convergence analysis is often omitted, since the algorithm performs very well in practice. Our convergence rates have several consequences for the pivoted Cholesky algorithm applied to matrices (see~\cref{sec:matrixcounterparts}), Gaussian process regression (see~\cref{sec:GaussianProcessRegression}), and the convergence of the power function in the P-greedy method (see~\cref{sec:Pgreedy}). We briefly discuss a few of these connections now. 

\subsection{Matrix counterparts of the kernel results}\label{sec:matrixcounterparts}

Our analysis has focused on kernels $K:\Omega\times\Omega\to\mathbb{R}$, where the pivoted Cholesky algorithm produces low-rank approximations with residuals controlled in terms of the fill distance of the pivot set. 
Interestingly, many of these ideas have natural discrete analogues when one replaces kernels by structured matrices. In this setting, the geometry of the domain $\Omega$ is mirrored by the index set $\{1,\dots,m\}$, and Lipschitz continuity of $K$ along the diagonal is replaced by discrete difference quotients of the matrix entries. To make this precise, let $A\in\mathbb{R}^{m\times m}$ be a SPD matrix, and define the discrete diagonal Lipschitz constant as
\begin{equation}
G_A = \max_{i\neq j}\,\frac{|A_{i,i}-A_{i,j}|}{|i-j|}.
\label{eq:GA}
\end{equation}
This quantity is the matrix analogue of the Lipschitz constant $L$ from the kernel setting: it measures how rapidly the entries of $A$ vary along each row as the column index changes. Applying the Cholesky algorithm with complete pivoting to $A$ then yields the following discrete counterpart of \Cref{thm:convergencegeneral}.

\begin{corollary}\label{cor:matrixversion}
    Let $A$ be an $m\times m$ SPD matrix and define $G_A$ as in~\cref{eq:GA}. Let $A_n$ be the rank-$n$ approximation obtained by the Cholesky decomposition with complete pivoting. 
    Then, for $1<n\geq m$, we have
    \[
        \|A - A_n\|_{\max} \;\leq\; \frac{4(m-1)G_A}{n-1} \;=\;\mathcal{O}(n^{-1}),
    \]
    where $\|\cdot\|_{\max}$ denotes the maximum absolute matrix entry.
\end{corollary}

This result shows that, in the matrix setting, complete pivoting necessarily drives the residual down at a uniform $\mathcal{O}(n^{-1})$ rate. 
Just as in the kernel case, the key mechanism is that complete pivoting prevents the pivots from clustering too closely: two pivot indices cannot be adjacent unless the residual is already small. In this way, the geometry of the index set enforces a quasi-uniform distribution of pivot indices, leading to algebraic decay of the residual. 

\subsection{Gaussian process regression}\label{sec:GaussianProcessRegression}

Gaussian process regression (Kriging)~\cite{williams2006gaussian} is a popular Bayesian nonparametric regression method, which is frequently used in machine learning. In the noise-free setting, one places a Gaussian process prior on the target function $f$ with mean $\mu:\Omega\to\mathbb{R}$ and covariance kernel $K:\Omega\times\Omega\to\mathbb{R}$. After $n$ observations, the posterior mean and covariance are updated via Bayes’ rule:
\[
\begin{aligned}
    \mu_n(x) &= \mu(x) + ([K(x, x_i)]_{i=1}^n)^T\left([K(x_i, x_j)]_{i, j=1}^n\right)^{-1}[y_i - \mu(x_i)]_{i=1}^n,\\
    R_n(x, x') &= K(x, x') - ([K(x, x_i)]_{i=1}^n)^T\left([K(x_i, x_j)]_{i, j=1}^n\right)^{-1}[K(x, x_i)]_{i=1}^n.
\end{aligned}
\]
Here, the posterior covariance $R_n$ coincides with the residual of the Cholesky algorithm after $n$ steps with pivot set $\{x_1,\dots,x_n\}$. Thus, bounds on the Cholesky residual directly translate into uniform variance reduction guarantees for Gaussian process regression.

Therefore, our convergence rates for the Cholesky algorithm with uniform pivoting, complete pivoting, or global/local maximum volume pivoting can be used to bound $R_n$ uniformly. In this setting, the control of the maximum variance in the Gaussian process is achieved after $n$ samples. 

Despite the empirical success of Gaussian process regression, its convergence properties are rarely addressed in the literature. The prevailing view is that smooth prior kernels yield satisfactory convergence in practice, see for instance~\cite{ambikasaran2015fast}. More systematic analyses connect Gaussian process regression to kernel interpolation, where the posterior mean corresponds to the interpolant and the posterior covariance coincides with the power function~\cite{scheuerer2013interpolation}. In this setting, deterministic convergence rates have been obtained for translation-invariant kernels~\cite{kanagawa2018gaussian,stuart2018posterior}, with few exceptions such as~\cite{de2005near}. Our results contribute to this picture by providing algebraic convergence rates for Gaussian process regression with general kernels, including those that are less smooth than $C^2$ and not necessarily translation-invariant. 

\subsection{The P-greedy algorithm in the kernel interpolation literature}\label{sec:Pgreedy}

The P-greedy algorithm is a widely used greedy strategy in kernel approximation and reduced basis methods~\cite{wenzel2023analysis}. Its selection rule is defined in terms of the so-called power function, which quantifies the worst-case interpolation error in the reproducing kernel Hilbert space (RKHS) associated with a kernel~\cite{de2005near,schaback2000adaptive}. 
We briefly recall these notions and then show that the P-greedy algorithm is equivalent to pivoted Cholesky.

Let $\Omega\subset \mathbb{R}^d$ and $K:\Omega\times\Omega\to\mathbb{R}$ be a continuous SPD kernel, with associated RKHS $\mathcal{H}_K$. 
Given a finite set of points $S_n=\{z_1,\dots,z_n\}\subset\Omega$, the kernel interpolant $I_{S_n} f\in \mathcal{H}_K$ to data $\{f(z_j)\}_{j=1}^n$ is the unique function of the form
\[
(I_{S_n} f)(x) \;=\; \sum_{j=1}^n \alpha_j K(x,z_j),
\]
that satisfies $(I_{S_n} f)(z_i)=f(z_i)$ for $i=1,\dots,n$.  
The error of this interpolant can be bounded in terms of the power function
\[
    P_{S_n}(x) \;=\; \sup_{\|f\|_{\mathcal{H}_K}=1}\, |f(x) - (I_{S_n} f)(x)|, \qquad x\in\Omega,
\]
which measures the maximal interpolation error at $x$ over all unit-norm functions in $\mathcal{H}_K$.

Let $K(S_n,S_n)=[K(z_i,z_j)]_{i,j=1}^n$, $k(x) = (K(x,z_1),\dots,K(x,z_n))^\top$, and define the Nystr\"{o}m approximation
\[
    \widehat K_{S_n}(x,y) \;=\; k(x)^\top K(S_n,S_n)^{-1}\, k(y).
\]
The residual kernel is the Schur complement
\begin{equation}\label{eq:residual-kernel}
    R_n(x,y) \;=\; K(x,y) - k(x)^\top K(S_n,S_n)^{-1} k(y).
\end{equation}
Note that if $S_n=\{z_1,\dots,z_n\}$ is chosen as the set of Cholesky pivots, then $R_n$ is exactly the Cholesky residual kernel after $n$ steps.

\begin{lemma}\label{lem:diag-equals-power2}
Let $\Omega\subset\mathbb{R}^d$ be compact and $K:\Omega\times \Omega \rightarrow \mathbb{R}$ a continuous SPD kernel. For every $x\in\Omega$ and every finite set of $n$ points $S_n\subset\Omega$, we have
\[
    R_n(x,x) = \big(P_{S_n}(x)\big)^2,
\]
where $R_n$ is the $n$th residual of the Cholesky algorithm that pivoted at $S_n$ and $P_{S_n}$ is the associated power function to the RKHS $\mathcal{H}_K$. 
\end{lemma}
\begin{proof}
Let $r_x(\cdot) = K(\cdot,x) - \sum_{j=1}^n \beta_j K(\cdot,z_j)$ with $\beta = K(S_n,S_n)^{-1} k(x)$.  
By the reproducing property of $\mathcal{H}_K$, for all $f\in \mathcal{H}_K$,
\[
f(x) - (I_{S_n} f)(x) = \langle f, r_x \rangle_{\mathcal{H}_K},
\]
where $\langle \cdot, \cdot \rangle_{\mathcal{H}_K}$ is the $\mathcal{H}_K$ inner-product. 
Hence, by Cauchy--Schwarz and optimality,
\[
P_{S_n}(x) = \|r_x\|_{\mathcal{H}_K}.
\]
Expanding $\|r_x\|_{\mathcal{H}_K}^2$ using the definition of $\beta$ gives
\[
\|r_x\|_{\mathcal{H}_K}^2 = K(x,x) - k(x)^\top K(S_n,S_n)^{-1} k(x).
\]
By~\cref{eq:residual-kernel}, this equals $R_n(x,x)$. Thus $P_{S_n}(x)^2 = R_n(x,x)$.
\end{proof}

By~\Cref{lem:diag-equals-power2}, the squared power function at $x$ coincides with the diagonal residual $R_n(x,x)$. Therefore, the P-greedy selection rule used in the kernel literature and given by
\[
    z_{n+1} = \arg\max_{x\in \Omega} P_{S_n}(x)
\]
is identical to the Cholesky complete pivoting rule
\[
    z_{n+1} = \arg\max_{x\in \Omega} R_n(x,x).
\]
In other words, P-greedy and the Cholesky with complete pivoting select the same pivots, and are essentially the same algorithm viewed through two different lenses: kernel approximation versus function approximation. 

Our convergence results can therefore be stated in terms of bounding the power function. In particular, for a kernel satisfying the assumptions in~\Cref{thm:convergencegeneral}, we know that
\[
    \|P_{S_n}\|_\infty^2 = \|R_n\|_\infty \;\leq\; \frac{8LR}{n^{1/d}-1} = \mathcal{O}(n^{-1/d}), 
\]
whenever the pivot set $S_n$ is generated by complete pivoting (equivalently, by P-greedy selection). This provides a deterministic, dimension-dependent rate for the decay of the power function under very weak smoothness assumptions.

Santin and Haasdonk~\cite{santin6convergence} analyze the P-greedy method for translation-invariant kernels, using connections to Kolmogorov $n$-widths and greedy approximation in Banach spaces. Their results yield near-optimal convergence rates depending on the smoothness of the kernel, and show that P-greedy points become asymptotically uniformly distributed in $\Omega$. Compared with their bounds, our results apply to more general (non-translation-invariant) kernels, and remain valid even under low regularity (e.g., Lipschitz continuity on the diagonal). On the other hand, their bounds are sharper in the high-smoothness, translation-invariant setting. Thus, the two approaches are complementary: our theory provides robust, dimension-dependent convergence guarantees under weak assumptions, while the P-greedy literature offers refined rates for smooth kernels in structured settings.

\section{Conclusion}\label{sec:conclusion}

We proved that the pivoted Cholesky algorithm converges under minimal smoothness assumptions. For SPD kernels that are Lipschitz on the diagonal, the residual is controlled by the fill distance of the pivots (see~\Cref{thm:filldistancebound}), yielding
\[
\|R_n\|_\infty=\mathcal{O}(n^{-1/d})
\]
for complete pivoting (see~\Cref{thm:convergencegeneral}). With $C^{1,1}$ regularity, the diagonal of the residual is locally quadratic in the distance to the nearest pivot (see~\Cref{lem:boundedregiondiff2}), giving the improved rate of
\[
\|R_n\|_\infty=\mathcal{O}(n^{-2/d})
\]
for both uniform and complete pivoting (see~\Cref{thm:filldistance-squared}). We also established a discrete analogue for SPD matrices with a diagonal Lipschitz constant $G_A$, obtaining $\|A-A_n\|_{\max}=\mathcal{O}(n^{-1})$ (see~\Cref{cor:matrixversion}).

Practically, approximate ($\delta$-complete) and grid-based pivoting retain the same algebraic rates up to constants, justifying standard implementations. The results transfer directly to Gaussian process regression (uniform variance reduction) and clarify the equivalence between complete pivoting and P-greedy via the power function. Open directions include rank-revealing refinements beyond uniform bounds, understanding staircase effects in convergence plots, and extending stability and noise-robustness guarantees.


\section*{Acknowledgments} 
A.T.~is grateful for support from NSF CAREER (DMS-2045646), DARPA-PA-24-04-07, and the Office of Naval Research under Grant Number N00014-23-1-2729. 

\bibliographystyle{siam}
\bibliography{refs}

\end{document}
